\section{Die These: in natürlichen Einheiten ist alpha = 1}

\textbf{[K]} Die stärkste Aussage über die Feinstrukturkonstante ist die
einfachste: in natürlichen Einheiten ist
\[
  \alpha = 1 .
\]
Der berühmte Wert $1/137{,}036$ ist \emph{kein} Naturgeheimnis, sondern die
Buchhaltung, die beim Übergang in das SI-System anfällt -- so wie
$c = 3\cdot10^8$ die Buchhaltung der Meter-Sekunden-Wahl ist.

\textbf{Das ist keine Erfindung dieser Theorie.} $\alpha$ auf Eins zu setzen ist
eine \emph{bekannte} Einheitenkonvention: man normiert die Ladung so, dass
$e^2 = 4\pi$ (rationalisierte Heaviside-Lorentz-Einheiten mit
$\varepsilon_0 = \mu_0 = 1$). In diesen Einheiten ist $\alpha = 1$ per
Konstruktion. Die Theorie übernimmt diese Konvention, sie begründet sie nicht neu.

\textbf{Der Grund ist aber ein anderer als bei $c, \hbar, k_B$.} Bei
$c, \hbar, k_B$ werden die
Umrechnungsfaktoren Eins, weil die Zeit-Masse-Dualität ihre Dimensionen
identifiziert (A085). $\alpha$ ist dimensionslos; diese Begründung greift hier
\emph{nicht}. Die Konvention $\alpha = 1$ ruht auf der elektromagnetischen
Ladungsnormierung -- ein eigener Grund, der nichts mit der Zeit-Masse-Dualität
zu tun hat und darum hier gesondert geführt wird, nicht beim Natureinheiten-Thema.

\section{Wie alpha = 1 hergeleitet wird}

\textbf{[B]} Die Herleitung läuft über die Normierung der Ladung und die
Maxwell-Dualität, nicht über eine Setzung. In SI ist
\[
  \alpha = \frac{e^2}{4\pi\varepsilon_0\hbar c}.
\]
Mit $\hbar = c = 1$ und der natürlichen Wahl $\varepsilon_0 = 1$ folgt aus
$\alpha = 1$ die Ladungsnormierung $e^2 = 4\pi$. Die Maxwell-Beziehung
$c^2 = 1/(\varepsilon_0\mu_0)$ liefert mit $c = 1$ und $\varepsilon_0 = 1$ dann
$\mu_0 = 1$. Einsetzen bestätigt:
\[
  \alpha = \frac{e^2}{4\pi\varepsilon_0\hbar c} = \frac{4\pi}{4\pi\cdot1\cdot1\cdot1} = 1 .
\]

\textbf{[K]} Der Kern ist die elektromagnetische Dualität
$\tfrac{1}{\varepsilon_0 c} = \mu_0 c$, die in natürlichen Einheiten perfekt
erfüllt ist ($1 = 1$). $\alpha = 1$ ist der Ausdruck dieser Balance zwischen
elektrischer ($\varepsilon_0$) und magnetischer ($\mu_0$) Kopplung.

\textbf{Die ehrliche Einordnung.} Der Schritt selbst ist Standard: die
Ladungsnormierung $e^2 = 4\pi$ ist die rationalisierte Konvention, in der
$\alpha = 1$ herauskommt. Was daran nicht Standard ist, ist die \emph{Deutung}:
dass der SI-Wert $1/137$ damit kein fundamentales Mysterium der Kopplung ist,
sondern die historische Ladungseinheit misst. Diese Deutung teilt nicht jeder --
die klassische Duff-Position hält $\alpha$ wegen seiner Dimensionslosigkeit für
absolut. Der eigentliche Beitrag der Theorie liegt aber nicht hier, sondern im
nächsten Abschnitt: dass der SI-Wert nicht nur Ladungseinheit misst, sondern über
$\alpha = \xi E_0^2$ den geometrischen Parameter $\xi$ trägt.

\textbf{Damit fällt die Mystifizierung weg.} Die Frage \enquote{warum gerade
$1/137$?} ist eine Frage an das SI-System, nicht an die Natur -- von derselben
Art wie \enquote{warum gerade $299\,792\,458$ Meter pro Sekunde?}. Der Zahlenwert
misst die historische Ladungseinheit, nicht eine Eigenschaft der Kopplung
selbst.

\section{Was übrig bleibt: xi und eine Skala}

\textbf{[K]} Wenn alle Konstanten in natürlichen Einheiten Eins sind, bleibt als
einziger echter Inhalt der eine dimensionslose Parameter $\xi$ und eine Skala.
Der SI-Wert von $\alpha$ trägt genau diese beiden Größen:
\[
  \alpha_\text{SI} = \xi\cdot\Bigl(\frac{E_0}{1\,\text{MeV}}\Bigr)^2 .
\]

Hier ist $\alpha$ nicht mehr die native Eins, sondern ihre SI-Kleidung. In dieser
Kleidung sieht man, dass $\alpha$ und $\xi$ dieselbe Größe kodieren, verbunden
durch die Skala $E_0$: $\xi$ ist die geometrische Lesart (Packungsverhältnis,
A020), $\alpha_\text{SI}$ die elektromagnetische. Zwei Seiten einer Medaille,
ein Freiheitsgrad -- so wie Masse, Energie, Temperatur und inverse Zeit eine
Größe in vier Kleidungen sind (A085).

\section{Warum der Messwert unwichtig ist}

\textbf{[K]} Die Theorie rechnet mit Verhältnissen (A080). Verhältnisse sind
einheitenfrei und hängen allein an der Geometrie. Der absolute Wert von $\alpha$
-- $1/137{,}036$ -- ist ein Absolutwert und gehört damit in die
Buchhaltungsschicht, nicht in den Kern.

\textbf{[K]} Die \emph{Formelzusammenhänge}, um die es geht, sind in SI-Einheiten
und im Standardmodell längst bekannt und bestätigt. Die Theorie erfindet sie
nicht, sie ordnet sie: $\alpha = \xi\,E_0^2$ verbindet bekannte Größen. Was sie
behauptet, ist die \emph{Struktur} dieser Verbindung -- dass $\alpha$ und $\xi$
derselbe Freiheitsgrad sind --, nicht einen neuen Zahlenwert für $\alpha$.

\textbf{[B]} Daraus folgt, wie die Zahlen der folgenden Abschnitte zu lesen sind:
\emph{nicht} als Test, den die Theorie besteht oder nicht besteht. Sie zeigen die
Genauigkeit, mit der die Buchhaltung den SI-Wert reproduziert -- eine Aussage
über die Brückenkonstante, nicht über die Struktur. Solange man mit Verhältnissen
rechnet, sind sie unwichtig.

\textbf{[SETZUNG]} Der Test der Theorie liegt anderswo. $\alpha$ ist keine der
vier Prognosegrößen der Serie (A010); diese sind Verhältnisse und
Strukturaussagen ($m_\tau$, die $\xi$-Ordnung der CHSH-Abweichung, $D_f$, $L_0$).
Ob $\alpha$ auf drei oder fünf Stellen an $1/137$ herankommt, prüft die Theorie
nicht -- es prüft die SI-Umrechnung.

\section{Warum die Zirkularität keine Schwäche ist}

\textbf{[SETZUNG]} Sobald man in der SI-Kleidung rechnet, ist das System
zirkulär, und das ist richtig so. Jedes System, das seine eigenen Konstanten
festlegt, ist zirkulär -- es fragt sich nur, wo man beginnt. (In natürlichen
Einheiten stellt sich die Frage gar nicht: dort ist $\alpha = 1$ und nichts zu
bestimmen.)

Man kann von $\xi$ ausgehen (geometrische Perspektive): dann ist $\xi$ primär und
$\alpha = \xi E_0^2$ abgeleitet. Man kann von $\alpha$ ausgehen
(elektromagnetische Perspektive): dann ist $\alpha$ primär und
$\xi = \alpha/E_0^2$ abgeleitet. Beide Wege liefern für alle anderen Konstanten
dieselben Werte. Sie sind mathematisch äquivalent.

\textbf{[B]} Die Frage \enquote{ist $\alpha$ ohne Zirkel aus $\xi$ ableitbar?}
ist deshalb falsch gestellt. In einem geschlossenen System gibt es kein Außen,
von dem her man ableiten könnte. Es gibt einen Eintrittspunkt, den man
deklariert, und den Rest, der daraus folgt. Die klassische Physik verbirgt das,
indem sie mehrere Konstanten als \enquote{fundamental} nebeneinanderstellt; diese
Theorie legt offen, dass es nur einen Eintrittspunkt braucht.

\textbf{Ein Fehler wäre hier trotzdem möglich,} nämlich den Eintrittspunkt für
eine Ableitung aus dem Nichts auszugeben. Wer $E_0 = \sqrt{\alpha/\xi}$ einsetzt
und daraufhin \enquote{exakt der gemessene Wert von $\alpha$} verkündet,
verwechselt das Ablesen der Konsistenz mit einem Beweis. Diesem Fehler entgeht
die $E_0$-Bestimmung, weil sie \emph{nicht} über $\sqrt{\alpha/\xi}$ läuft:
$E_0$ folgt aus den Leptonmassen und dem aus $\xi$ hergeleiteten $K_\text{frak}$
(folgender Abschnitt), und $\alpha$ wird damit auf $0{,}017\,\%$ vorwärts
getroffen. Die Zirkularität ist an dieser Stelle also dadurch gebrochen, dass
$K_\text{frak}$ nicht aus $\alpha$ stammt. \textit{(Korrigiert am 30.~Sept.~2026: vorher \enquote{der exakte $E_0$ ist von Weg~2 unabhängig von $\alpha$ getroffen, und die Übereinstimmung der beiden $\alpha$-freien Wege macht die $\alpha$-Stelle überbestimmt} -- das beruhte auf der Formel $E_0^2 = 4\sqrt2\,m_\mu/\xi^{p}$, die entfernt ist, siehe folgender Abschnitt.)}

\section{Die Skala E0 als Eintrittspunkt}

\textbf{[K]} Was $\xi$ von $\alpha$ trennt, ist allein die Skala $E_0$. Für sie
gibt es zwei eng verwandte Wege, die beide $\alpha$ nicht benutzen: dasselbe
geometrische Mittel der Leptonmassen, einmal ohne und einmal mit fraktaler
Korrektur.

\textbf{[K]} \emph{Weg~1 (ohne Korrektur).} Das geometrische Mittel der beiden
leichtesten Leptonmassen,
\[
  E_0^{(1)} = \sqrt{m_e\,m_\mu} = 7{,}348\ \text{MeV},
\]
liefert eingesetzt $\alpha = \xi\cdot(7{,}348)^2 = 1/138{,}9$, also $-1{,}35\,\%$
neben dem Messwert. Das ist die Variante, in der sich $K_\text{frak}$
herauskürzt: als Verhältnisrechnung trägt sie den Faktor nicht (Dok.~122).

\textbf{[K]} \emph{Weg~2 (mit fraktaler Korrektur).} Mit dem in A040 aus $\xi$
hergeleiteten Faktor $K_\text{frak} = 1-100\xi = 74/75$,
\[
  E_0^{(2)} = \sqrt{\frac{m_e\,m_\mu}{K_\text{frak}}}
  = \frac{7{,}348}{\sqrt{0{,}98667}} = 7{,}397\ \text{MeV},
\]
folgt $\alpha^{-1} = K_\text{frak}/(\xi\,m_e m_\mu) = 137{,}06$, also $+0{,}017\,\%$
neben dem Messwert; mit dem gerundeten Anker $E_0 = 7{,}398$~MeV ergibt sich
$137{,}035$. Die Wurzel steht, weil $E_0$ quadratisch in $\alpha$ eingeht:
$K_\text{frak}$ wirkt auf $\alpha^{-1}$. \textit{(Korrigiert am 30.~Sept.~2026: Weg~2 war vorher \enquote{eine Konstruktion aus $\xi$ und $m_\mu$ allein}, $E_0^2 = 4\sqrt2\,m_\mu/\xi^{p}$, die $7{,}398$~MeV \enquote{ohne $m_e$ und ohne Rückrechnung aus $\alpha$} treffe. Diese Formel ist als zweiter Weg entfernt: $p$ ist nirgends hergeleitet, das Prüfskript setzte $p = -0{,}2679$, das genau $7{,}398$ trifft; mit $p = 4$ ergäbe sie $E_0 = 1{,}4\times10^{9}$~MeV, und sie ist dimensional unstimmig (Dok.~011, 041). Im übrigen Korpus findet sich keine Herleitung. Weg~2 ist das fraktal korrigierte Mittel (Korpus R72).)}

\section{Die zwei Wege legen $K_\text{frak}$ fest}

\textbf{[K]} Die beiden Wege unterscheiden sich nur um $K_\text{frak}$. Stellt man
die Frage umgekehrt -- welcher Faktor muss zwischen ihnen stehen, damit $\alpha$
exakt getroffen wird? --, dann ist die Antwort eindeutig:
\[
  K_\text{frak} \;=\; \xi\,m_e m_\mu\,\alpha^{-1}
  \;=\; \Bigl(\frac{E_0^{\,(1)}}{\sqrt{\alpha/\xi}}\Bigr)^{2}
  \;=\; \Bigl(\frac{7{,}3479}{7{,}3980}\Bigr)^{2} \;=\; 0{,}98650 .
\]

\textbf{[K]} Dieser Wert trifft den in A040 unabhängig hergeleiteten Faktor
\[
  K_\text{frak} = 1 - 100\xi = \frac{74}{75} = 0{,}98667
\]
auf $1{,}7\cdot10^{-4}$; anders gelesen verlangt $\alpha$ vor $\xi$ den
Koeffizienten $101{,}2$ statt $100$. Die beiden Wege sind damit der
indirekte Beweis, dass $K_\text{frak}$ nur diesen Wert haben kann: A040 gewinnt ihn aus dem kumulativen RG-Lauf und der
Konsistenzbedingung an die effektive fraktale Dimension, der $\alpha$-Sektor
verlangt denselben Wert als Lücke zwischen dem nackten Mittel und dem aus
$\alpha$ folgenden $E_0$.

\textbf{[B]} Damit hat $K_\text{frak}$ im $\alpha$-Sektor einen \emph{zweiten
Zeugen}. Diese Bestimmung benutzt $\alpha$ -- sie ist keine $\alpha$-freie
Messung --, aber sie benutzt $K_\text{frak}$ nicht; was in A040 eine
hergeleitete Konstante ist, wird hier aus $\alpha$ und den Leptonmassen
\emph{abgelesen}. \textit{(Korrigiert am 30.~Sept.~2026: vorher hießen die beiden Wege \enquote{voneinander unabhängig} und die $\alpha$-Stelle \enquote{überbestimmt}; $K_\text{frak} = (7{,}398/7{,}348)^{-2} = 0{,}9862$ galt als $\alpha$-freie \enquote{Messung}, die sich mit $1-100\xi$ auf $8\cdot10^{-5}$ treffe, und dieser Rest als \enquote{Test der Polmasse}. Das beruhte auf der entfernten Weg-2-Formel. $7{,}398 = \sqrt{\alpha/\xi}$ enthält $\alpha$; richtig ist $(7{,}3479/7{,}3980)^2 = 0{,}98650$, also $1{,}7\cdot10^{-4}$ neben $74/75$.)}

\section{Eine Koinzidenz, der man widerstehen muss}

\textbf{[K]} Es liegt nahe, nach einem \emph{schöneren} Ausdruck für $E_0$ zu
suchen als dem Massenweg, und es gibt einen auffälligen: die Euler-Zahl im
Quadrat,
\[
  \mathrm{e}^2 = 7{,}389\ \text{MeV}\ (?),\qquad
  \alpha = \xi\cdot(7{,}389)^2 = \frac{1}{137{,}37}\quad(-0{,}24\,\%),
\]
der den SI-Wert sogar besser trifft als der Massenweg. Diesem Ausdruck ist zu
\emph{widerstehen}, und zwar aus einem prinzipiellen Grund: er setzt eine reine
Zahl (die Euler-Zahl) mit einer Energie in der willkürlichen Einheit MeV gleich.
Ein MeV ist eine menschliche Festlegung; keine reine Zahl kann eine physikalische
Skala \emph{in} MeV festlegen, ohne dass die Übereinstimmung von der Einheitenwahl
abhinge. In einer anderen Energieeinheit stünde eine andere Zahl, und die
Koinzidenz verschwände.

\textbf{[K]} Das ist genau das Muster, vor dem der Korpus an anderer Stelle warnt
(A105): eine Zahl, rückwärts an einen schönen Ausdruck angepasst, ohne
Vorwärtsgehalt. Dass $\mathrm{e}^2$ näher trifft als $\sqrt{m_e m_\mu}$, macht es
\emph{nicht} zum besseren Kandidaten -- der Massenweg trägt einen physikalischen
Grund (die Leptonmassen), die Euler-Koinzidenz trägt keinen. Sie wird hier nur
festgehalten, um sie ausdrücklich \emph{nicht} als Verankerung zu buchen. Der
schöne Ausdruck wird zudem gar nicht gebraucht: $E_0$ ist bereits über
Weg~2 (Massenmittel mit $K_\text{frak}$) mit Vorwärtsgehalt bestimmt, während
$\mathrm{e}^2$ nur eine Zahl an eine Einheit hält.

\section{Prüfskript}

\begin{itemize}
  \item Die Äquivalenz beider Perspektiven ($\alpha = \xi E_0^2$ und
    $\xi = \alpha/E_0^2$ liefern identische Werte); der $\alpha$-freie
    harmonische Weg mit $\kappa = 4$; die zwei Zahlenwege; und die ausdrückliche
    Prüfung, dass die Einsetzung von $\sqrt{\alpha/\xi}$ eine Konsistenz und kein
    unabhängiger Beweis ist:
    \texttt{python/A\_Serie\_Skripte/a130\_alpha\_kette.py}
\end{itemize}

\section{Was offen bleibt}

\textbf{Erstens ist $E_0$ $\alpha$-frei bestimmt, aber nicht überbestimmt.}
Weg~2 liefert $E_0 = 7{,}397$~MeV aus den Leptonmassen und dem aus $\xi$
hergeleiteten $K_\text{frak}$, ohne $\alpha$; damit folgt $\alpha$ auf
$0{,}017\,\%$. Offen bleibt der Rest von $1{,}7\cdot10^{-4}$ zwischen dem von
$\alpha$ verlangten $K_\text{frak} = 0{,}98650$ und $74/75$. Die
Euler-Koinzidenz $E_0 = \mathrm{e}^2$ bleibt ausdrücklich \emph{keine} Herleitung
und wird auch nicht gebraucht. \textit{(Korrigiert am 30.~Sept.~2026: vorher \enquote{Der exakte $E_0 = 7{,}398$ wird von Weg~2 aus $\xi$ und $m_\mu$ getroffen, ohne $\alpha$ und ohne $m_e$ \ldots\ überbestimmt} -- beruhte auf der entfernten Formel $E_0^2 = 4\sqrt2\,m_\mu/\xi^{p}$.)}

\textbf{Zweitens bleibt der absolute Anker an gemessene Massen gebunden.}
Weg~2 fixiert $E_0$ aus $m_e$, $m_\mu$ und $\xi$; die absolute Skala tritt also
über die gemessenen Leptonmassen ein. Eine Bestimmung des exakten $E_0$ aus
rein dimensionsloser Geometrie ohne jeden Massenanker gibt es nicht -- das ist
aber die allgemeine Struktur der SI-Brücke (A080), kein Defizit speziell des
$\alpha$-Sektors. Da die Theorie mit Verhältnissen prüft, ist der verbleibende
Rest für ihre Bewertung ohne Gewicht.

\textbf{Drittens gilt eine Untergrenze.} Die relative Beschreibung
$\alpha \leftrightarrow \xi$ ist bis zur Sub-Planck-Grenze $L_0 = \xi\,\ell_P$
gesichert (A210); darunter macht die Theorie keine Aussage.

\section{Ersetzt}

Dieses Dokument nimmt auf: Dok.~011 (Feinstrukturkonstante), Dok.~043 (Auflösung
der Konstanten), Dok.~044, Dok.~101 (Zirkularität der Konstanten, Kernthese: ein
Freiheitsgrad, $\xi$ und $\alpha$ äquivalent). Eingearbeitet: P40.

\section{Verwendung von KI-Werkzeugen}

Dieses Dokument ist unter Verwendung KI-gestützter Werkzeuge entstanden. Die
Fragestellungen, die Setzungen und alle inhaltlichen Entscheidungen stammen vom
Autor; die Werkzeuge formulieren aus, rechnen nach und erstellen Prüfskripte.
Die Verantwortung für Inhalt und Fehler liegt vollständig beim Autor. Der
ausführliche Wortlaut steht in A010.
